% Provenance (never printed): sources/book2/BRIEFS.md, g12-12. The 1960s relation; the claim of a stable trade-off and its
% refutation in the 1970s; adaptive expectations and the vertical long-run curve (proved by a recurrence); disinflation and
% the sacrifice ratio; anchored expectations and 2021-25. Real numbers: figdata/grade-12/12-phillips-curve/ (BLS CPI-U
% annual inflation, the measure of chapter 10; JST unemployment to 1986; WDI unemployment 2015-2025). Ledger: Phillips 1958
% (bibliographic, as the name of the relation), Volcker's appointment (Federal Reserve History), Book 3's row on Ball 1994.
% The natural rate of 6 % in the problem is given, not estimated.
\chapter{The Phillips Curve and Its Breakdown}\label{ch:g12:phillips-curve}

In 1961 American unemployment was 6.7\,\% and consumer prices rose by 1.1\,\%. By 1969 unemployment had fallen to 3.5\,\% and inflation had risen to 5.5\,\%. For a decade,
lower unemployment came with higher inflation, year after year, almost along a line. In the next decade the line broke: in 1975 unemployment was 8.5\,\% and inflation
9.1\,\%; in 1980, 7.1\,\% and 13.5\,\%.
% figdata: phil6069 (1961 6.7/1.1; 1969 3.5/5.5); phil7083 (1975 8.5/9.1; 1980 7.1/13.5)

This chapter follows a relation that seemed to offer a choice, the evidence that withdrew the offer, and the model that explained why.

\begin{openingchart}{United States: unemployment rate against consumer price inflation, per cent, each year joined to the next; 1960--1969 (blue) and 1970--1983 (red).}{U.S.
Bureau of Labor Statistics (CPI-U); Jordà--Schularick--Taylor Macrohistory Database (unemployment)}
\begin{tikzpicture}
\begin{axis}[omscatter, width=0.72\linewidth, height=6.4cm, xmin=3, xmax=10.5, ymin=0, ymax=15, xlabel={unemployment rate, per cent},
  ylabel={inflation, per cent}, xtick={3,...,10}, ytick={0,5,10,15}]
\addplot[omDef, thick, sharp plot, mark=*, mark size=1.3pt] table[x=u, y=pi] {figdata/grade-12/12-phillips-curve/out/phil6069.csv};
\addplot[omThm, thick, sharp plot, mark=*, mark size=1.3pt] table[x=u, y=pi] {figdata/grade-12/12-phillips-curve/out/phil7083.csv};
\node[font=\scriptsize, omDef, anchor=north] at (axis cs:6.7,1.1) {1961};
\node[font=\scriptsize, omDef, anchor=south east] at (axis cs:3.5,5.5) {1969};
\node[font=\scriptsize, omThm, anchor=west] at (axis cs:8.5,9.1) {1975};
\node[font=\scriptsize, omThm, anchor=south] at (axis cs:7.1,13.5) {1980};
\node[font=\scriptsize, omThm, anchor=west] at (axis cs:9.6,3.2) {1983};
\end{axis}
\end{tikzpicture}
\end{openingchart}

\begin{recall}
From \cref{ch:g11:unemployment}: the \omterm{def:g11:unemployment:natural-rate}{natural rate of unemployment}. From \cref{ch:g12:aggregate-demand-and-supply}: demand and supply shocks, and the return to the long
run. From \cref{ch:g12:monetary-policy}: the Federal Reserve's policy in 1979--83 and 2021--23.
\end{recall}

\section{A relation in the data}\label{sec:g12:phillips-curve:relation}

\begin{definition}[Phillips curve]\label{def:g12:phillips-curve:phillips-curve}
A \emph{Phillips curve}\index{Phillips curve} is a relation in which inflation is higher when unemployment is lower. The name comes from a 1958 study of British wages and
unemployment from 1861 to 1957.
\end{definition}
% ledger: g12:phillips1958

The reason is the one of \cref{ch:g12:aggregate-demand-and-supply}: when demand is strong, firms hire, unemployment falls, workers can ask for more, and firms raise prices
faster. The blue years of the opening chart fit it well.

\begin{model}{The stable Phillips trade-off}
\assumes One fixed relation between inflation and unemployment.
\predicts A menu: a government that accepts more inflation can have permanently lower unemployment.
\tested Refuted. From 1970 to 1983 the points left the 1960s line: in 1975 and 1980 unemployment and inflation were both high (the red years of the opening chart).
\breaks Whenever people come to expect inflation: the relation shifts up.
\end{model}

\begin{definition}[Stagflation]\label{def:g12:phillips-curve:stagflation}
\emph{Stagflation}\index{stagflation} is high inflation together with high unemployment and weak growth.
\end{definition}

\begin{record}[Stagflation, 1973--82]
% figdata: phil7083 (1974 5.6/11.1; 1975 8.5/9.1; 1980 7.1/13.5; 1982 9.7/6.1)
After the oil embargo of 1973 (\cref{ch:g12:aggregate-demand-and-supply}), American inflation reached 11.1\,\% in 1974, and unemployment 8.5\,\% in 1975. Inflation never
returned to the levels of the early 1960s: it was 7.6\,\% in 1978 and 13.5\,\% in 1980, with unemployment at 7.1\,\%. Two supply shocks and inflation that people had learned
to expect lifted the whole relation.
\end{record}

\section{Expectations}\label{sec:g12:phillips-curve:expectations}

\begin{definition}[Adaptive expectations]\label{def:g12:phillips-curve:adaptive}
People have \emph{adaptive expectations}\index{adaptive expectations} when they expect inflation this year to equal last year's: $\pi^e_t = \pi_{t-1}$.
\end{definition}

\begin{definition}[Expectations-augmented Phillips curve]\label{def:g12:phillips-curve:augmented}
The \emph{expectations-augmented Phillips curve}\index{expectations-augmented Phillips curve} is the relation $\pi_t = \pi^e_t - \beta\,(u_t - u^n) + s_t$, where $u^n$ is the
\omterm{def:g11:unemployment:natural-rate}{natural rate of unemployment}, $\beta > 0$, and $s_t$ is a supply shock (zero in normal times).
\end{definition}

\begin{notation}[$\beta$ in this chapter]\label{not:g12:phillips-curve:beta}
Elsewhere in the series $\beta$ is a discount factor or a regression coefficient; in this chapter it is the slope of the \omterm{def:g12:phillips-curve:phillips-curve}{Phillips curve}.
\end{notation}

Wage setters who expect 5\,\% inflation build it into their wage claims; inflation then exceeds 5\,\% only if unemployment is below the natural rate.

\begin{proposition}[No long-run trade-off]\label{prop:g12:phillips-curve:no-trade-off}
With \omterm{def:g12:phillips-curve:adaptive}{adaptive expectations} and no supply shocks, if unemployment is held at $u^n - x$ with $x > 0$, inflation rises by $\beta x$ every year: $\pi_t = \pi_0 + t\beta x$, without
limit. Inflation is constant only when $u = u^n$.
\end{proposition}
\begin{proof}
With $\pi^e_t = \pi_{t-1}$ and $u_t - u^n = -x$, the curve gives $\pi_t = \pi_{t-1} + \beta x$, an arithmetic sequence: $\pi_t = \pi_0 + t\beta x$, which grows without
limit since $\beta x > 0$. Inflation is constant, $\pi_t = \pi_{t-1}$, exactly when $\beta(u_t - u^n) = 0$, that is $u_t = u^n$.
\end{proof}

\begin{omfigure}
\begin{tikzpicture}
\begin{groupplot}[group style={group size=2 by 1, horizontal sep=1.2cm}, width=0.47\linewidth, height=5.2cm, xmin=2, xmax=9, ymin=0, ymax=9,
  xlabel={unemployment}, tick label style={font=\scriptsize}, label style={font=\scriptsize}, xtick={4,5,7}, ytick={2,3,4,5,6}]
\nextgroupplot[title={\scriptsize holding unemployment at 4\,\%}, ylabel={inflation}]
\addplot[gray, thick] coordinates {(5,0) (5,9)};
\addplot[omDef, domain=2:9, samples=2] {2 - 0.5*(x - 5)};
\addplot[omDef, dashed, domain=2:9, samples=2] {2.5 - 0.5*(x - 5)};
\addplot[omDef, dotted, thick, domain=2:9, samples=2] {3 - 0.5*(x - 5)};
\addplot[omThm, only marks, mark=*, mark size=1.6pt] coordinates {(5,2) (4,2.5) (4,3) (4,3.5)};
\node[font=\tiny, anchor=west] at (axis cs:5,8.4) {long run};
\nextgroupplot[title={\scriptsize a disinflation}]
\addplot[gray, thick] coordinates {(5,0) (5,9)};
\addplot[omDef, domain=2:9, samples=2] {6 - 0.5*(x - 5)};
\addplot[omDef, dashed, domain=2:9, samples=2] {4 - 0.5*(x - 5)};
\addplot[omThm, thick, mark=*, mark size=1.6pt] coordinates {(5,6) (7,5) (7,4) (7,3) (5,3)};
\end{groupplot}
\end{tikzpicture}

{\small Short-run curves $\pi = \pi^e - 0.5(u - 5)$ with $u^n = 5$. Left: holding unemployment at 4\,\% raises inflation from 2 to 2.5, 3 and 3.5\,\% as expectations catch up;
each year's curve lies above the last. Right: from 6\,\% inflation, holding unemployment at 7\,\% for three years brings inflation down to 3\,\%; unemployment then returns to
5\,\% at the lower inflation.}
\end{omfigure}

\begin{model}{The expectations-augmented Phillips curve}
\assumes \omterm{def:g12:phillips-curve:adaptive}{Adaptive expectations} and a \omterm{def:g11:unemployment:natural-rate}{natural rate of unemployment}.
\predicts No lasting trade-off: unemployment held below the natural rate brings ever-rising inflation; bringing inflation down requires a period of unemployment above it.
\tested The 1970s, when inflation kept rising with unemployment not low; 1980--86, when inflation fell after unemployment rose far above its earlier levels (the weekend
  problem).
\breaks When expectations no longer follow last year's inflation but stay near a target, the curve stops shifting with last year's inflation, and in
  the 2010s it also looked flat; supply shocks shift it, as in 2021--22.
\end{model}

\section{The cost of bringing inflation down}\label{sec:g12:phillips-curve:disinflation}

\begin{definition}[Disinflation and sacrifice ratio]\label{def:g12:phillips-curve:disinflation}
A \emph{disinflation}\index{disinflation} is a lasting fall in the rate of inflation. The \emph{sacrifice ratio}\index{sacrifice ratio}%
\phantomsection\label{def:g12:phillips-curve:sacrifice-ratio} is the cumulative unemployment above the natural rate during a disinflation, in point-years, divided by the fall
in inflation, in points.
\end{definition}

\begin{record}[The disinflation of 1980--86 and its cost]
% ledger: b1:fed-targets, b1:ball1994; figdata: volcker, phillips_key (excess_8086 13.7, fall 11.6, ratio 1.18)
On 6 October 1979 the Federal Reserve changed its operating procedure to bring inflation down. Inflation fell from 13.5\,\% in 1980 to 3.2\,\% in 1983 and 1.9\,\% in 1986; unemployment, 5.8\,\% in
1979, reached 9.7\,\% in 1982 and was still 7.0\,\% in 1986. With a natural rate of 6\,\%, unemployment above it adds up to 13.7 point-years, for a fall of 11.6 points of
inflation: a \omterm{def:g12:phillips-curve:disinflation}{sacrifice ratio} of about 1.2. Across 65 \omterm{def:g12:phillips-curve:disinflation}{disinflations} in rich countries with moderate inflation, the ratio (measured there as lost output) was smaller when inflation was brought down
quickly and when wages were set more flexibly.
\end{record}

\section{Anchored expectations}\label{sec:g12:phillips-curve:anchored}

\begin{definition}[Anchored expectations]\label{def:g12:phillips-curve:anchored}
Expectations are \emph{anchored}\index{anchored expectations} when people expect inflation to stay near the central bank's target whatever inflation was last year.
\end{definition}

With $\pi^e$ fixed at the target, the augmented curve no longer accelerates: low unemployment raises inflation a little, for as long as it lasts, and no more. In the 2010s
American unemployment fell from 5.3\,\% in 2015 to 3.7\,\% in 2019 while inflation stayed near 2\,\%.

\begin{omchart}{United States: unemployment rate against consumer price inflation, per cent, 2015--2025, each year joined to the next; inflation in 2025 is from the eleven monthly index values
published.}{World Bank, World Development Indicators
(SL.UEM.TOTL.NE.ZS); U.S. Bureau of Labor Statistics (CPI-U)}
\begin{tikzpicture}
\begin{axis}[omscatter, width=0.62\linewidth, height=5.6cm, xmin=3, xmax=9, ymin=0, ymax=9, xlabel={unemployment rate, per cent},
  ylabel={inflation, per cent}, xtick={3,...,9}, ytick={0,2,...,8}]
\addplot[omDef, thick, sharp plot, mark=*, mark size=1.3pt] table[x=u, y=pi] {figdata/grade-12/12-phillips-curve/out/phil1525.csv};
\node[font=\scriptsize, anchor=west] at (axis cs:5.4,0.3) {2015};
\node[font=\scriptsize, anchor=north east] at (axis cs:3.7,1.8) {2019};
\node[font=\scriptsize, anchor=west] at (axis cs:8.1,1.2) {2020};
\node[font=\scriptsize, anchor=east] at (axis cs:3.6,8) {2022};
\node[font=\scriptsize, anchor=west] at (axis cs:4.3,2.6) {2025};
\end{axis}
\end{tikzpicture}
\end{omchart}

\begin{record}[2021--25]
% figdata: phil1525 (2021 5.3/4.7; 2022 3.6/8.0; 2023 3.6/4.1; 2024 4.0/2.9; 2025 4.3/2.6)
Inflation rose to 4.7\,\% in 2021 and 8.0\,\% in 2022 as unemployment fell to 3.6\,\%, with supply shocks pushing the curve up. It then fell to 4.1\,\% in 2023, 2.9\,\% in 2024
and 2.6\,\% in 2025 (from the eleven monthly values published), while unemployment rose only to 4.3\,\%. A fall of five points of inflation with so little unemployment does not fit the \omterm{def:g12:phillips-curve:disinflation}{sacrifice ratio} of the 1980s;
it fits a supply shock that faded while expectations stayed close to the target. Which of the two mattered more is still argued over.
\end{record}

\section{Exercises}

\begin{exercise}[$\star$]\label{exo:g12:phillips-curve:1}
From the opening chart, give unemployment and inflation in 1961 and in 1969.
\end{exercise}

\begin{exercise}[$\star$]\label{exo:g12:phillips-curve:2}
From the opening chart, compare 1980 with 1965 (unemployment 4.5\,\%, inflation 1.6\,\%). Does 1980 lie on the 1960s line?
\end{exercise}

\begin{exercise}[$\star$]\label{exo:g12:phillips-curve:3}
The American consumer price index averaged 270.97 in 2021 and 292.65 in 2022. Find inflation in 2022.
\end{exercise}

\begin{exercise}[$\star$]\label{exo:g12:phillips-curve:4}
The curve is $\pi = 3 - 0.5(u - 5)$. At what unemployment rate is inflation equal to the expected 3\,\%? What is inflation when $u = 4$?
\end{exercise}

\begin{exercise}[$\star$]\label{exo:g12:phillips-curve:5}
During a \omterm{def:g12:phillips-curve:disinflation}{disinflation}, unemployment is above the natural rate by 2, 3 and 1 points in three years, and inflation falls by 4 points. Find the \omterm{def:g12:phillips-curve:disinflation}{sacrifice ratio}.
\end{exercise}

\begin{exercise}[$\star\star$]\label{exo:g12:phillips-curve:6}
With $\beta = 0.5$, $u^n = 5$, \omterm{def:g12:phillips-curve:adaptive}{adaptive expectations} and inflation of 2\,\% last year, unemployment is held at 4\,\%. Find inflation in the next three years.
\end{exercise}

\begin{exercise}[$\star\star$]\label{exo:g12:phillips-curve:7}
With $\beta = 0.5$, $u^n = 5$ and inflation of 5\,\% last year, unemployment is held at 6\,\% for three years. Find inflation each year.
\end{exercise}

\begin{exercise}[$\star\star$]\label{exo:g12:phillips-curve:8}
After \cref{exo:g12:phillips-curve:6}, unemployment returns to 5\,\%. What happens to inflation? What was gained, and what was the lasting cost?
\end{exercise}

\begin{exercise}[$\star\star$]\label{exo:g12:phillips-curve:9}
Using the weekend problem's table, compute the \omterm{def:g12:phillips-curve:disinflation}{sacrifice ratio} for 1980--83 only, with a natural rate of 6\,\%.
\end{exercise}

\begin{exercise}[$\star\star$]\label{exo:g12:phillips-curve:10}
From the chart of 2015--2025, compare 2019 and 2022. Did the same unemployment rate come with the same inflation?
\end{exercise}

\begin{exercise}[$\star\star$]\label{exo:g12:phillips-curve:11}
Redo \cref{exo:g12:phillips-curve:6} with expectations \omterm{def:g12:phillips-curve:anchored}{anchored} at 2\,\%.
\end{exercise}

\begin{exercise}[$\star\star\star$]\label{exo:g12:phillips-curve:12}
State the claim of a stable trade-off, and the evidence that contradicts it.
\end{exercise}

\begin{exercise}[$\star\star\star$]\label{exo:g12:phillips-curve:13}
With $\beta = 0.5$, $u^n = 5$, \omterm{def:g12:phillips-curve:adaptive}{adaptive expectations} and inflation of 2\,\%, a supply shock $s = 3$ hits in one year, with unemployment at 5\,\%. Find inflation that year and
the next, if unemployment stays at 5\,\%. What would bring it back to 2\,\%?
\end{exercise}

\begin{exercise}[$\star\star\star$]\label{exo:g12:phillips-curve:14}
In the proposition, what happens to $\pi_t$ as $t$ grows? Why does this mean that no unemployment rate below $u^n$ can be held with stable inflation?
\end{exercise}

\begin{exercise}[$\star\star\star$]\label{exo:g12:phillips-curve:15}
Give two reasons why the \omterm{def:g12:phillips-curve:phillips-curve}{Phillips curve} of the 2010s looked flat.
\end{exercise}

\section{Problem: The Disinflation of 1979--1986}

\begin{problem}[{Weekend problem --- the cost of bringing inflation down}]
\label{pb:g12:phillips-curve:1}
The table gives American unemployment and consumer price inflation, and unemployment above a natural rate of 6\,\%, taken as given (all in per cent).

\medskip
\begin{center}
\small
\pgfplotstabletypeset[col sep=comma, columns={year,u,pi,excess},
  columns/year/.style={column name=year, int detect, column type=l},
  columns/u/.style={column name=unemployment, fixed, precision=1, zerofill},
  columns/pi/.style={column name=inflation, fixed, precision=1, zerofill},
  columns/excess/.style={column name={above 6\,\%}, fixed, precision=1, zerofill},
  every head row/.style={before row=\toprule, after row=\midrule}, every last row/.style={after row=\bottomrule}]
  {figdata/grade-12/12-phillips-curve/out/volcker.csv}
\end{center}

\medskip
\noindent\textbf{Part I --- The record.}
\begin{enumerate}
\item In which year was inflation highest? Unemployment?
\item By how many points did inflation fall from 1980 to 1986?
\item By how many from 1980 to 1983?
\item What was inflation in 1979, the year of the Federal Reserve's change of operating procedure?
\item Was unemployment back at 6\,\% by 1986?
\end{enumerate}

\medskip
\noindent\textbf{Part II --- The sacrifice ratio.}
\begin{enumerate}[resume]
\item Add up unemployment above 6\,\% from 1980 to 1986.
\item Find the \omterm{def:g12:phillips-curve:disinflation}{sacrifice ratio}.
\item Redo it with a natural rate of 5.5\,\%.
\item Redo it with a natural rate of 6.5\,\%.
\item What does this say about the precision of a \omterm{def:g12:phillips-curve:disinflation}{sacrifice ratio}?
\end{enumerate}

\medskip
\noindent\textbf{Part III --- The curve.}
\begin{enumerate}[resume]
\item With \omterm{def:g12:phillips-curve:adaptive}{adaptive expectations} and $u^n = 6$, use 1982 to estimate $\beta$ from $\pi_{1982} - \pi_{1981} = -\beta(u_{1982} - 6)$.
\item With that $\beta$, predict inflation in 1983 from 1982, and compare.
\item Predict inflation in 1984 from 1983 the same way, and compare. What went wrong?
\item How would a credible announcement of lower inflation change expectations, and the \omterm{def:g12:phillips-curve:disinflation}{sacrifice ratio}?
\item Was this \omterm{def:g12:phillips-curve:disinflation}{disinflation} quick or gradual? What does the cross-country evidence say about the cost of each?
\end{enumerate}

\medskip
\noindent\textbf{Part IV --- Interpretation.}
\begin{enumerate}[resume]
\item Who bore the cost of the \omterm{def:g12:phillips-curve:disinflation}{disinflation}?
\item Why might the cost have been larger without the earlier change in the Federal Reserve's procedures and its stated aim?
\item From 2022 to 2025 inflation fell by 5.4 points while unemployment never rose above 4.3\,\%. Why can the 1980s ratio not be applied?
\item Which model of this chapter does this problem test, and does it pass?
\item State the result of the problem: the \omterm{def:g12:phillips-curve:disinflation}{sacrifice ratio} of the American \omterm{def:g12:phillips-curve:disinflation}{disinflation}, in point-years of unemployment per point of inflation.
\end{enumerate}
\end{problem}
